% Generated draft: requires mathematical and source review.
\documentclass[11pt]{article}
\usepackage[margin=25mm]{geometry}
\usepackage{fontspec}
\setmainfont{Latin Modern Roman}
\usepackage{amsmath,amssymb,mathtools,unicode-math}
\setmathfont{Latin Modern Math}
\usepackage{enumitem,needspace,xurl,hyperref}
\setlength{\parindent}{0pt}
\setlength{\parskip}{5pt}
\setlength{\emergencystretch}{3em}
\newcommand{\N}{\mathbb{N}}
\newcommand{\Z}{\mathbb{Z}}
\newcommand{\Q}{\mathbb{Q}}
\newcommand{\R}{\mathbb{R}}
\newcommand{\C}{\mathbb{C}}
\newcommand{\F}{\mathbb{F}}
\newcommand{\A}{\mathbb{A}}
\newcommand{\PP}{\mathbb{P}}
\newcommand{\E}{\mathbb{E}}
\newcommand{\eps}{\varepsilon}
\newcommand{\dd}{\,\mathrm{d}}
\begin{document}
\section*{UnsolvedMath ID: 30006292}
\textbf{Source problem code:} catalogue entry
\subsection{Definitions and mathematical statement}
The curated question concerns the usual two-dimensional incompressible Navier--Stokes equations. Their differential expression is
\[
\partial_t u+(u\cdot\nabla)u+\nabla p=\nu\Delta u+f,
\qquad \nabla\cdot u=0,
\]
where \(u\) is the velocity, \(p\) is the pressure, and \(\nu>0\) is the viscosity. The spatial domain, boundary conditions, forcing \(f\), initial-data space, and associated phase space are not specified in the supplied record. This display therefore describes the equations only; it does not select an initial-boundary-value problem.

For a specified autonomous realization generating a solution semigroup \(S(t):H\to H\), \(t\geq0\), on a normed phase space \(H\), a common semigroup formulation defines an \emph{inertial manifold} to be a nonempty finite-dimensional Lipschitz submanifold \(\mathcal M\subset H\) that is positively invariant,
\[
S(t)\mathcal M\subseteq\mathcal M\qquad(t\geq0),
\]
and exponentially attracts every bounded subset of \(H\): there exists \(\beta>0\) such that for every nonempty bounded \(B\subset H\) there exists \(C_B<\infty\) satisfying
\[
\sup_{u_0\in B}\inf_{v\in\mathcal M}
\|S(t)u_0-v\|_H\leq C_B e^{-\beta t}
\qquad(t\geq0).
\]
This is an explicit working convention, not a definition supplied by the source. Whether the intended definition additionally requires a global graph over specified modes, stronger regularity, or an exponential tracking property remains unresolved.

The source question is: \emph{Does the usual two-dimensional Navier--Stokes evolution admit an inertial manifold?} Under the working convention above, for the intended realization this asks whether there exist \(\mathcal M\) and \(\beta>0\) with the stated properties. The supplied record does not identify that realization or quantify over possible domains, boundary conditions, viscosities, and forcings. Consequently, it does not determine a unique fully specified existence assertion. In particular, no universal quantifier over all such choices is asserted here. An autonomous semigroup formulation also requires a suitable forcing specification, which is absent.

The extracted original passage refers to the Navier--Stokes equations without explicitly specifying dimension; the curated statement specifies two dimensions. The supplied status is \emph{open}. Its supplied literature assessment distinguishes the usual equations from regularized or hyperviscous models, for which it reports known inertial-manifold results. These are source qualifications and supplied assessments, not independently verified mathematical or literature conclusions.
\subsection{Short English statement}
The curated problem asks whether the usual two-dimensional Navier--Stokes equations possess a finite-dimensional invariant manifold that attracts solutions exponentially. The domain, boundary conditions, forcing, phase space, and precise inertial-manifold convention are missing, so the intended existence question still requires mathematical clarification. The supplied original excerpt does not explicitly state the dimension.
\subsection{Sources}
\begin{itemize}
\item Local registry: \texttt{data/unsolvedmath/problems.json}, permanent UnsolvedMath ID \(30006292\), source code \texttt{catalogue entry}, title \emph{Inertial Manifolds for Navier--Stokes Equations}; collection: Oberwolfach Reports Open Problems. The supplied record labels the problem open and labels its statement \texttt{corrected\_verified}, with the qualification that only formula delimiters, notation, grammar, and typography were repaired. These labels are supplied registry metadata.
\item \emph{Mathematical Advances in Geophysical Fluid Dynamics} (2025), Oberwolfach Reports, pp.\ 1215--1270, DOI 10.4171/owr/2025/24. The extracted passage is attributed to the contribution \emph{Inertial Manifolds for Regularized Navier--Stokes Equations}, Yanqiu Guo, joint work with Mohammad Abu Hamed, Ciprian G. Gal, Michael Ilyin, and Edriss S. Titi, p.\ 1255. Report link, also supplied as the literature source ``original OWR report'': \url{https://doi.org/10.4171/owr/2025/24}. The registry supplies proposed year 2025 but no separate proposer attribution.
\item Supplied literature source: ``Zelik, attractors and inertial manifolds,'' \url{https://www.mathnet.ru/eng/rm10095}.
\item Literature assessment date: \texttt{2026-08-22}. The supplied dated curation labels the problem open and reports that inertial manifolds are known for regularized or hyperviscous Navier--Stokes models, while existence for the usual two-dimensional equations remains a long-standing open problem. The background describes this as dated literature triage derived from the supplied OWR curation and requests independent verification of time-sensitive status and cited claims. This assessment and these references have not been independently verified here.
\end{itemize}
\end{document}
